% TOP_PROBLEM: {"problemId": "problem.hardy-littlewood-prime-k-tuple-conjecture", "canonicalTitle": "Hardy-Littlewood Prime k-Tuple Conjecture", "releaseRank": 18, "primaryDomain": "number_theory_arithmetic_geometry", "primaryDomainLabel": "Number theory & arithmetic geometry", "displayStatus": "open", "releaseStatus": "open"}
% !TeX program = lualatex
% ATTEMPT_STATUS: unresolved
% Research continuation by GPT_6_Astra_Ultra, 27 September 2026.
% Catalogue statements and existing sources preserved verbatim.
\documentclass[11pt]{article}
\usepackage[margin=25mm]{geometry}
\usepackage{fontspec}
\setmainfont{Latin Modern Roman}
\usepackage{amsmath,amssymb,mathtools}
\usepackage{unicode-math}
\setmathfont{Latin Modern Math}
\usepackage{xurl,hyperref}
\hypersetup{colorlinks=true,urlcolor=blue}
\setcounter{secnumdepth}{0}
\setlength{\parindent}{0pt}
\setlength{\parskip}{5pt}
\setlength{\emergencystretch}{3em}
\begin{document}
\section*{\texorpdfstring{Hardy-Littlewood Prime k-Tuple Conjecture}{Hardy-Littlewood Prime k-Tuple Conjecture}}\label{18}
\subsection{Definitions and mathematical statement}
For a finite set $H=\{h_1,\ldots,h_k\}\subset{\mathbb{Z}}$ of distinct offsets, let $\nu_p(H)$ be the number of residues it occupies modulo $p$. Admissibility means $\nu_p(H)<p$ for every prime $p$. Define
\[
 \mathfrak S(H)=\prod_p\frac{1-\nu_p(H)/p}{(1-1/p)^k},\qquad
 \pi_H(x)=\#\{1\le n\le x:n+h_i\text{ prime for every }i\}.
\]
For each fixed admissible $H$, the conjecture is $\pi_H(x)\sim\mathfrak S(H)x/(\log x)^k$ as $x\to\infty$. Here $f\sim g$ means $f/g\to1$; this is stronger than infinitude alone.
\par\smallskip\textit{Formulation and concept references:} \href{https://mathworld.wolfram.com/k-TupleConjecture.html}{[S1]}.

\subsection{Short English statement}
Does every locally admissible fixed pattern of primes occur with exactly the frequency predicted by its divisibility corrections?
\subsection{Research attempt}

\paragraph{Exact target and literature distinction.}
The offsets are distinct and fixed before $x$ tends to infinity. Negative
offsets merely remove finitely many initial integers: put
$n_0=\max(1,2-\min H)$ so all $n+h$ under consideration are at least two.
The empty pattern, if permitted, gives $\pi_H(x)=\lfloor x\rfloor$ and
is immediate. We assume $k\ge1$ below. Repeated offsets cannot be inserted
as extra coordinates: two copies of zero would leave the prime count
unchanged while incorrectly changing the proposed logarithmic exponent.

Maynard [E1] proves strong results about many primes in bounded tuples
and a positive proportion of patterns satisfying an infinitude statement.
Those conclusions do not assert the precise asymptotic for each prescribed
$H$ requested here. No primorial-counting claim in catalogue source [S3]
is assumed to prove that stronger assertion. The selected attack is to
derive everything supplied by local divisibility exactly, then attempt
to pass to the sieve level that detects actual primes. This isolates a
quantifiable missing term, rather than treating local independence as an
unproved probability law.

\paragraph{Admissibility and positivity of the singular series.}
If $\nu_p(H)=p$ for some prime $p$, every integer $n$ makes some
$n+h_i$ divisible by $p$. If all these integers are prime, that particular
coordinate must equal $p$, so $n=p-h_i$ for one of at most $k$ choices.
Thus an inadmissible pattern has only finitely many occurrences. This
proves the necessity of local admissibility for infinitude, allowing
exceptional small tuples such as $(3,5,7)$ for $H=\{0,2,4\}$.

Conversely, assume admissibility. Let
$B=\max_{i,j}|h_i-h_j|$, taking $B=0$ for $k=1$, and choose
$y>\max(B,2k,2)$. For $p>y$ all offsets are distinct modulo $p$, so
$\nu_p(H)=k$. Write
$F_p=(1-k/p)/(1-1/p)^k$. Expanding the logarithms absolutely gives
\[
 \log F_p=-\sum_{m\ge2}\frac{k^m-k}{mp^m},\qquad
 -\frac{k^2}{p^2}\le\log F_p\le0.
\]
Indeed $1/m\le1/2$ for $m\ge2$ and $k/p<1/2$, giving
$\sum_{m\ge2}(k/p)^m/m\le k^2/p^2$. Hence the infinite product
converges to a positive value, and its tail has the explicit enclosure
\begin{equation}\label{ra:18:tail}
 \exp\!\left(-\frac{k^2}{y-1}\right)
 \le\prod_{p>y}F_p\le1.
\end{equation}
We used the elementary bound
$\sum_{p>y}p^{-2}\le\sum_{n>y}n^{-2}\le1/(y-1)$.
All earlier factors are positive by admissibility. This proves that the
catalogue's constant is well-defined and nonzero, with a computable
tail certificate. It proves no occurrence of the pattern.

Translation of $H$ preserves every $\nu_p$ and the singular series.
Multiplying all offsets by an integer need not preserve it: additional
collisions arise at primes dividing that integer. The proposed frequency
therefore contains arithmetic information beyond the diameter of the
pattern. In the special case $k=1$ every local factor is exactly one.

\paragraph{A finite admissibility test is not a finite primality test.}
Admissibility itself is decidable by checking only primes $p\le k$:
for $p>k$ the elementary inequality $\nu_p(H)\le k<p$ is automatic.
For example, $H=\{0,2,6\}$ occupies one residue modulo two and two
residues modulo three, so it is admissible without any further local
tests. The pattern $\{0,2,4\}$ occupies all three residues modulo three
and fails. This finite procedure checks whether a local obstruction
exists, not whether the admissible pattern has a single prime occurrence
beyond its initial range. There is no analogous cutoff at $p=k$ when
testing primality of the values $n+h$: composite values can have their
smallest prime factor much larger than the number of coordinates.
The distinction matters computationally as well. A program may certify
admissibility quickly and compute a positive singular-series enclosure,
yet still lack any lower bound for the occurrence count. Those certificates
are useful input checks, and the next calculation states exactly how far
their counting consequences extend.

\paragraph{Finite wheels: what the Chinese remainder theorem proves.}
Let $P(z)=\prod_{p\le z}p$ and
\[
 T_H(x,z)=\#\{n_0\le n\le x:
                   \gcd(\prod_{h\in H}(n+h),P(z))=1\}.
\]
For each prime $p\le z$ there are exactly $p-\nu_p(H)$ allowed
residues for $n$. The Chinese remainder theorem identifies these choices
independently across the distinct prime moduli. Thus a complete block
modulo $P(z)$ has exactly
\[
 R_H(z)=\prod_{p\le z}(p-\nu_p(H))
\]
survivors. In an interval of length $L=\lfloor x\rfloor-n_0+1\ge0$,
each allowed residue appears either $\lfloor L/P(z)\rfloor$ or
$\lceil L/P(z)\rceil$ times. Therefore
\begin{equation}\label{ra:18:wheel}
 T_H(x,z)=L D_H(z)+O(R_H(z)),\qquad
 D_H(z)=\prod_{p\le z}(1-\nu_p(H)/p),
\end{equation}
with absolute error at most $R_H(z)$.

For every fixed $z$, this is an exact density theorem as $x\to\infty$.
Its implied error depends on $z$; the statement does not license taking
$z=\sqrt x$ while retaining a bounded error. Admissibility ensures that
every finite wheel has survivors, but most survivors can still have prime
factors beyond the wheel. The infinite sequence of positive finite-wheel
densities is not itself a prime-pattern count.

\paragraph{Exact inclusion--exclusion exposes the accumulated error.}
For squarefree $d\mid P(z)$ put
$\nu_d=\prod_{p\mid d}\nu_p(H)$, with $\nu_1=1$, and let
$A_d(x)$ count $n_0\le n\le x$ such that $d\mid\prod_h(n+h)$.
CRT now counts forbidden residues, giving
$A_d=L\nu_d/d+e_d$ with $|e_d|\le\nu_d$.
Inclusion--exclusion gives the exact identities
\begin{align}
 T_H(x,z)&=\sum_{d\mid P(z)}\mu(d)A_d(x)
           =L D_H(z)+\mathcal R_H(x,z),\label{ra:18:IE}\\
 \mathcal R_H(x,z)&=\sum_{d\mid P(z)}\mu(d)e_d,
 \qquad |\mathcal R_H(x,z)|\le\prod_{p\le z}(1+\nu_p(H)).\notag
\end{align}
The expansion is finite, so there is no convergence issue. The difficulty
is the size and signed behavior of the remainder when $z$ grows. A bound
on each progression separately does not survive summing over the enormous
family of divisors at the needed accuracy.

Now set $z=\sqrt{x+\max H}$ for sufficiently large $x$. Every positive
integer $n+h\le x+\max H$ which is composite has a prime divisor at
most $z$. Hence every survivor is an actual prime tuple, with all its
components greater than $z$. Conversely any prime tuple omitted by the
wheel has one component at most $z$, hence $n\le z-\min H$.
We have proved the unconditional relation
\begin{equation}\label{ra:18:prime-level}
 \pi_H(x)=T_H(x,\sqrt{x+\max H})+O_H(\sqrt x).
\end{equation}
The error is negligible compared with $x/(\log x)^k$ for fixed $k$.
This brings the argument to the exact sieve level where the desired
asymptotic would follow, but the remainder in \eqref{ra:18:IE} is now
uncontrolled.

\paragraph{A sharper falsification of simply discarding that remainder.}
Mertens' classical product theorem is
$\prod_{p\le z}(1-1/p)\sim e^{-\gamma}/\log z$; see [E2], introduction.
Together with the proved convergence of the singular series, it gives
\[
 D_H(z)\sim\mathfrak S(H)\frac{e^{-k\gamma}}{(\log z)^k}.
\]
At the prime-detecting level $z=\sqrt{x+\max H}$ this predicts
\[
 L D_H(z)\sim (2e^{-\gamma})^k\mathfrak S(H)
                          \frac{x}{(\log x)^k}.
\]
The multiplier is not one. Already for $H=\{0\}$, dropping
$\mathcal R_H$ contradicts the prime number theorem [E3]: the actual
count is asymptotic to $x/\log x$, whereas $2e^{-\gamma}=1.1229\ldots$.
Thus the issue is stronger than a missing small-error proof. At this
sieve level the signed correction has a nonzero main term even in the
solved one-coordinate case.

For general fixed admissible $H$, the precise estimate that would complete
this particular calculation is
\begin{equation}\label{ra:18:gap}
 \mathcal R_H(x,\sqrt{x+\max H})
 =\bigl(1-(2e^{-\gamma})^k+o(1)\bigr)
                  \mathfrak S(H)\frac{x}{(\log x)^k}.
\end{equation}
Equations \eqref{ra:18:IE} and \eqref{ra:18:prime-level} prove its
sufficiency. No derivation of \eqref{ra:18:gap} is given. It is a
specific arithmetic correction problem, not a consequence of multiplying
the finite local densities. This elementary obstruction is consistent
with the limitations of sieve methods discussed in [E1].

\paragraph{Stress tests for order of limits and pattern choice.}
An abstract family makes the order-of-limits issue unmistakable. Let
$f_z(x)=1$ when $x\ge z^3$ and zero otherwise. For each fixed $z$,
$f_z(x)\to1$ as $x\to\infty$, but $f_{\sqrt x}(x)=0$ for $x>1$.
Thus a theorem valid for each fixed primorial does not become a theorem
at a primorial whose cutoff grows with the observation range. The actual
wheel has additional arithmetic structure, but that structure must enter
a uniform estimate; the quantifiers alone do not provide it.

One also cannot choose a new favorable pattern $H_x$ at each scale and
claim the required result for every fixed $H$. The bounds
\eqref{ra:18:tail} and \eqref{ra:18:prime-level} explicitly depend on
the offsets and their diameter. Conversely, a proof for one prescribed
admissible pattern with $k\ge2$ would be genuine progress, even without
uniformity in all offsets. The catalogue asks for all such fixed patterns,
not a density-one family of them and not merely an unspecified pair of
prime positions inside a larger tuple.

\paragraph{Finite verification.}
The standard-library script checks CRT and complete inclusion--exclusion
for the wheel $P(7)=210$, for $H=\{0\},\{0,2\},\{0,2,6\},
\{0,2,4\},\{-2,0,4\}$ and $x=100,1000,10000$.
It also independently sieves the actual prime tuples. Selected results are
\[
\begin{array}{c|r|r|r|r}
H&R_H(7)&x&\pi_H(x)&T_H(x,7)\\\hline
\{0,2\}&15&100&8&6\\
\{0,2\}&15&1000&35&70\\
\{0,2\}&15&10000&205&713\\
\{0,2,6\}&8&10000&55&381\\
\{0,2,4\}&0&10000&1&0
\end{array}
\]
At small $x$ a wheel can exclude genuine tuples containing its small
primes, explaining the first row. At larger $x$ it admits many composites,
explaining the later discrepancies. The inadmissible example retains its
one exceptional tuple. All counts and identities are exact; no asymptotic
constant was fitted to the data. Reproduction files are in
\texttt{\_Local/top\_problems/continuation\_15\_25\_2026\_09\_27/analytic/}.

\paragraph{Remaining gap and outcome.}
The proved results are positivity with a quantitative tail bound for the
singular series, exact finite-wheel counts, and the prime-level sieve
identity. They isolate the missing correction \eqref{ra:18:gap} and
show why setting it to zero gives the wrong answer even for one prime.
Neither the required signed remainder nor a counterexample pattern has
been obtained. The next mathematical advance would need uniform control
of that remainder or a different method detecting joint primality,
while retaining the precise leading constant.

\textbf{Outcome: UNRESOLVED.} The local computations and obstruction
proofs are rigorous partial analysis, not a solution of the asymptotic
conjecture. Direct derivation, primary references, and exact finite checks
were used; no independent formal certification is claimed.

\subsection{Sources}
\begin{itemize}
\item[S1]k-Tuple Conjecture (MathWorld).
\url{https://mathworld.wolfram.com/k-TupleConjecture.html}.
\textit{Formulation source.}
\item[S2]Correlations of almost primes — Institut Henri Poincaré seminar abstract.
\url{https://www.ihp.fr/en/node/3096}.
\textit{Further reference; not a proof endorsement.}
\item[S3]Proof of the Hardy-Littlewood K-tuple Conjecture in the Distribution of Numbers Coprime with the Primorial — Tim Samshuijzen, May 2026 revision.
\url{https://rxiv.org/pdf/2501.0157v4.pdf}.
\textit{Further reference; not a proof endorsement.}
\item[S4]Higher uniformity of arithmetic functions in short intervals II. Almost all intervals — Matomäki, Radziwiłł, Shao, Tao and Teräväinen (2026).
\url{https://link.springer.com/article/10.1007/s00222-026-01408-6}.
\textit{Further reference; not a proof endorsement.}

\item[E1]James Maynard, \emph{Small gaps between primes},
Annals of Mathematics 181 (2015), 383--413.
\url{https://annals.math.princeton.edu/wp-content/uploads/annals-v181-n1-p07-p.pdf}.
\textit{Primary research context: bounded tuples and infinitude for a positive
proportion of patterns, not the asymptotic for every specified pattern.
Consulted 27 September 2026.}
\item[E2]Youness Lamzouri, \emph{A bias in Mertens' product formula},
arXiv:1410.3777v2 (6 February 2015), introduction.
\url{https://arxiv.org/pdf/1410.3777v2}.
\textit{Research source recording the classical unconditional Mertens product
theorem used in the cutoff stress test. The paper's conditional bias results
are not invoked. Consulted 27 September 2026.}
\item[E3]Andrew V. Sutherland, \emph{18.785 Number Theory I, Lecture 16:
Riemann's zeta function and the prime number theorem}, MIT, fall 2021.
\url{https://math.mit.edu/classes/18.785/2021fa/LectureNotes16.pdf}.
\textit{Author lecture notes for the one-coordinate prime number theorem,
used to falsify the naive sieve asymptotic. Consulted 27 September 2026.}
\end{itemize}
\end{document}
